% !TeX program = xelatex
% Portable practice of book structure; not SJTUThesis.
\documentclass[UTF8,fontset=fandol,openany]{ctexbook}
\usepackage{amsmath,amsthm,graphicx}
\usepackage[hidelinks]{hyperref}
\newtheorem{theorem}{零误差性质}[chapter]
\title{温度校正排版练习}\author{实验者}\date{}
\begin{document}
\frontmatter
\maketitle
\chapter{摘要}本练习以六条合成观测演示学位论文结构。\chapter{Abstract}This exercise demonstrates thesis structure using six synthetic observations.
\tableofcontents\listoffigures\listoftables
\mainmatter
\chapter{方法}\label{chap:method}
\begin{theorem}若所有预测值等于真实值，则 RMSE 为零。\end{theorem}
\begin{proof}各平方误差均为零，代入定义即可。\end{proof}
\appendix
\chapter{数据说明}六条观测均为教学合成数据，不对应任何真实设备或受试者。
\backmatter
\chapter{致谢}感谢教学材料作者。
\end{document}
